% Solucionario, Tema 1 — Ondas electromagnéticas (Saleh 5.1, 5.2, 5.4, 2.2). Problemas 1 a 7.
% Convención de Saleh en todo el tema: E(r,t)=Re{E(r)e^{jωt}}, I=|E|^2/2η. El Prob. 4 compara con Boyd.

\noindent\textbf{Convención del tema.} Se usa la de Saleh y Teich: campos reales en caligráfica,
$\Ec(\rv,t)=\mathrm{Re}\{\E(\rv)e^{j\omega t}\}$, una onda que avanza en $+z$ es $e^{-jkz}$ e
$I=|E|^2/2\eta$ con $\eta=\eta_0/n$ (Saleh, ecs.~5.3-1 y 5.4-8, pp.~162 y 166). Solo en el
Problema~4 aparece además la de Boyd, $\tilde E=Ee^{-i\omega t}+\cc$ e $I=2n\epsilon_0c|E|^2$
(Boyd, ec.~2.2.16, p.~77), y allí se dice explícitamente.

%=====================================================================
\sol{1}{Ecuación de onda con fuente no lineal}{Saleh §5.1 (pp.~153--154) y §5.2B, ecs.~5.2-16 a 5.2-25 (pp.~159--161); Boyd §2.1, ecs.~2.1.9b--2.1.11 (pp.~71--72).}\label{sol:1}

\inc{a} \textbf{Deducción.} Partimos de las ecuaciones de Maxwell en un medio sin cargas ni
corrientes libres (Saleh, ecs.~5.1-7 a 5.1-10, p.~154):
\begin{equation*}
\nabla\times\Hc=\frac{\partial\Dc}{\partial t},\qquad
\nabla\times\Ec=-\frac{\partial\Bc}{\partial t},\qquad
\nabla\cdot\Dc=0,\qquad \nabla\cdot\Bc=0,
\end{equation*}
con las relaciones $\Dc=\epsilon_0\Ec+\Pc$ (5.1-11) y, por ser el medio no magnético
($\boldsymbol{\mathcal M}=0$), $\Bc=\mu_0\Hc$ (5.1-12).

\emph{Paso 1: rotacional de la ley de Faraday.}
$\nabla\times(\nabla\times\Ec)=-\nabla\times\partial_t\Bc=-\mu_0\,\partial_t(\nabla\times\Hc)$, donde se
intercambió $\nabla$ con $\partial_t$ (las coordenadas espaciales y el tiempo son independientes) y
se sacó $\mu_0$ porque es constante.

\emph{Paso 2: ley de Ampère--Maxwell.} Sustituyendo $\nabla\times\Hc=\partial_t\Dc$:
\begin{equation*}
\nabla\times(\nabla\times\Ec)=-\mu_0\frac{\partial^2\Dc}{\partial t^2}
=-\mu_0\epsilon_0\frac{\partial^2\Ec}{\partial t^2}-\mu_0\frac{\partial^2\Pc}{\partial t^2}.
\end{equation*}
Esta forma (Boyd, ec.~2.1.9b, p.~71) es completamente general: no supone nada del medio salvo que
no es magnético.

\emph{Paso 3: identidad vectorial.} Con $\nabla\times(\nabla\times\Ec)=\nabla(\nabla\cdot\Ec)-\nabla^2\Ec$
(Saleh, p.~153) se obtiene la ec.~5.2-24 (Saleh, p.~161):
\begin{equation*}
\nabla(\nabla\cdot\Ec)-\nabla^2\Ec=-\frac{1}{c_0^2}\frac{\partial^2\Ec}{\partial t^2}-\mu_0\frac{\partial^2\Pc}{\partial t^2},
\qquad \epsilon_0\mu_0=\frac1{c_0^2}.
\end{equation*}

\emph{Paso 4: aquí se usa que el medio es homogéneo e isótropo.} Hay que eliminar
$\nabla(\nabla\cdot\Ec)$. La única información sobre $\nabla\cdot\Ec$ es $\nabla\cdot\Dc=0$. En un medio
lineal $\Dc=\epsilon\Ec$; entonces
\begin{equation*}
0=\nabla\cdot(\epsilon\Ec)=\epsilon\,\nabla\cdot\Ec+\nabla\epsilon\cdot\Ec .
\end{equation*}
Para concluir $\nabla\cdot\Ec=0$ se necesitan \emph{dos} cosas: que $\epsilon$ sea un escalar (isotropía;
si fuera un tensor, $\nabla\cdot(\boldsymbol\epsilon\Ec)$ no se factoriza así) y que no dependa de la
posición, $\nabla\epsilon=0$ (homogeneidad). Con ambas, $\nabla\cdot\Ec=0$ y el término
$\nabla(\nabla\cdot\Ec)$ desaparece. Cambiando el signo de toda la ecuación:
\resultado{$\displaystyle \nabla^2\Ec-\frac{1}{c_0^2}\frac{\partial^2\Ec}{\partial t^2}=\mu_0\frac{\partial^2\Pc}{\partial t^2}$
\quad (Saleh, ec.~5.2-25, p.~161). El paso que exige homogeneidad e isotropía es
$\nabla\cdot\Dc=0\Rightarrow\nabla\cdot\Ec=0$.}
\emph{Sutileza.} En un medio \emph{no lineal}, aunque sea homogéneo e isótropo,
$\nabla\cdot\Ec=-\nabla\cdot\Pc_{NL}/\epsilon$ no es exactamente cero (Boyd, p.~71). Se desprecia
porque se anula para una onda plana transversal infinita y es pequeño cuando la envolvente varía
lentamente (Boyd, p.~72). Saleh lo da por cero sin comentarlo.

\inc{b} Escribimos $\Pc=\epsilon_0\chi^{(1)}\Ec+\Pc_{NL}$ (el primer término es la respuesta lineal,
no dispersiva, con $\chi^{(1)}$ constante) y sustituimos en el lado derecho de (5.2-25):
\begin{equation*}
\nabla^2\Ec-\frac{1}{c_0^2}\frac{\partial^2\Ec}{\partial t^2}
=\mu_0\epsilon_0\chi^{(1)}\frac{\partial^2\Ec}{\partial t^2}+\mu_0\frac{\partial^2\Pc_{NL}}{\partial t^2}
=\frac{\chi^{(1)}}{c_0^2}\frac{\partial^2\Ec}{\partial t^2}+\mu_0\frac{\partial^2\Pc_{NL}}{\partial t^2}.
\end{equation*}
Pasando el término lineal a la izquierda y agrupando:
\resultado{$\displaystyle \nabla^2\Ec-\frac{n^2}{c_0^2}\frac{\partial^2\Ec}{\partial t^2}=\mu_0\frac{\partial^2\Pc_{NL}}{\partial t^2},
\qquad n^2=1+\chi^{(1)}$ \quad (es la ec.~1.1.5 de Boyd, p.~4, con $\mu_0=1/\epsilon_0c^2$).}
\textbf{Interpretación.} El lado izquierdo es el operador de onda del medio lineal, con velocidad
$c_0/n$: la parte lineal de la polarización solo cambia la velocidad de la luz. El lado derecho es
un término fuente, igual que la corriente de una antena: $\Pc_{NL}$ es una densidad de dipolos y
$\partial_t\Pc_{NL}$ una densidad de corriente de polarización; su segunda derivada es la
\emph{aceleración} de las cargas ligadas, y una carga acelerada radia. Si $\Pc_{NL}$ oscila a
$2\omega$, $\omega_1+\omega_2$, etc., esos dipolos radian campos a esas frecuencias nuevas, que no
estaban en la onda incidente. Toda la óptica no lineal (SHG, SFG, efecto Kerr) consiste en resolver
esta ecuación con la fuente adecuada.

\inc{c} En un medio anisótropo lineal, $D_i=\sum_j\epsilon_{ij}E_j$ (Saleh, ec.~5.2-22, p.~160). Entonces
\begin{equation*}
\nabla\cdot\Dc=\sum_{i,j}\epsilon_{ij}\frac{\partial\mathcal E_j}{\partial x_i}=0,
\qquad\text{mientras que}\qquad
\nabla\cdot\Ec=\sum_{i}\frac{\partial\mathcal E_i}{\partial x_i}.
\end{equation*}
Son combinaciones distintas de derivadas: que la primera se anule no obliga a que lo haga la
segunda (salvo que $\epsilon_{ij}=\epsilon\,\delta_{ij}$). Para una onda plana $e^{-j\kv\cdot\rv}$:
$\nabla\cdot\Dc=0\Leftrightarrow\kv\cdot\D=0$, pero $\D$ y $\E$ no son paralelos, así que en general
$\kv\cdot\E\neq0$: $\E$ tiene una componente longitudinal (a lo largo de $\kv$). Es el origen del
\emph{walk-off}: el vector de Poynting $\E\times\Hv$ forma un ángulo $\rho$ con $\kv$.

En la práctica se \emph{desprecia el término $\nabla(\nabla\cdot\Ec)$}. Su tamaño relativo es
$|\nabla(\nabla\cdot\Ec)|/|\nabla^2\Ec|\sim|\kv\cdot\E|/(k|\E|)=\tan\rho$, de modo que la aproximación
vale cuando la anisotropía es débil, $|n_e-n_o|\ll n$, que es lo que hace pequeño el ángulo de walk-off
($\rho\lesssim$ unos grados en los cristales usados en ONL), o cuando la onda viaja a lo largo de un
eje principal ($\rho=0$, el término se anula exactamente). Boyd añade que también es pequeño
bajo la aproximación de envolvente lentamente variable (p.~72).

\inc{d} Para un medio inhomogéneo con $\epsilon(\rv)$ que varía lentamente frente a $\lambda$, el término
$\nabla[(1/\epsilon)\nabla\epsilon\cdot\Ec]$ de (5.2-18) es despreciable y queda (Saleh, ec.~5.2-19, p.~159)
\begin{equation*}
\nabla^2\Ec-\frac{n^2(\rv)}{c_0^2}\frac{\partial^2\Ec}{\partial t^2}\approx0 .
\end{equation*}
Con $n(\rv)=n+\Delta n(\rv)$ y $|\Delta n|\ll n$:
$n^2(\rv)=n^2+2n\Delta n+(\Delta n)^2\approx n^2+2n\Delta n$. Sustituyendo y pasando el término en
$\Delta n$ a la derecha:
\begin{equation*}
\nabla^2\Ec-\frac{n^2}{c_0^2}\frac{\partial^2\Ec}{\partial t^2}=\frac{2n\Delta n}{c_0^2}\frac{\partial^2\Ec}{\partial t^2}
=\mu_0\frac{\partial^2}{\partial t^2}\big(2\epsilon_0 n\Delta n\,\Ec\big),
\end{equation*}
donde se usó $1/c_0^2=\mu_0\epsilon_0$ y que $\Delta n$ no depende del tiempo. Comparando con la
forma de (b) ($\mu_0\partial_t^2\Pc_{NL}$):
\resultado{$\Delta\Pc=2\epsilon_0n\,\Delta n\,\Ec$, que es la ec.~5.2-20 de Saleh (p.~159). Otra forma de
verlo: $\Pc=\epsilon_0(n^2-1)\Ec\Rightarrow\Delta\Pc=\epsilon_0\Delta(n^2)\Ec=2\epsilon_0 n\Delta n\,\Ec$.}
Si $\Delta n$ depende de la intensidad ($\Delta n=n_2I$, efecto Kerr), esta misma fuente es una
polarización no lineal de tercer orden: es la puerta de entrada al cap.~4 de Boyd.

\idea{(5.2-25) vale para cualquier medio homogéneo, isótropo y no magnético, lineal o no; la
condición se usa \emph{solo} para tirar $\nabla(\nabla\cdot\Ec)$. Cuidado con los prefactores:
$\mu_0=1/(\epsilon_0c_0^2)$, así que la fuente se escribe igual en Saleh ($\mu_0\partial_t^2\Pc$) que
en Boyd ($\frac{1}{\epsilon_0c^2}\partial_t^2\tilde P$). Una perturbación del índice $\Delta n$ equivale
a $\Delta\Pc=2\epsilon_0n\Delta n\Ec$ (factor 2 por derivar $n^2$).}

%=====================================================================
\sol{2}{Clasificación de medios}{Saleh §5.2, recuadro de definiciones (p.~156), §5.2B (pp.~158--162) y oscilador de Lorentz, ecs.~5.5-15 a 5.5-19 (pp.~176--177); Boyd ec.~1.1.2 (p.~2).}\label{sol:2}

Criterios (Saleh, p.~156): \emph{lineal} si vale la superposición entre $\Pc$ y $\Ec$;
\emph{no dispersivo} si $\Pc(t)$ depende solo de $\Ec$ en el mismo instante (sin derivadas ni
integrales temporales); \emph{espacialmente no dispersivo} si $\Pc(\rv)$ depende solo de $\Ec$ en el mismo
punto (sin derivadas espaciales); \emph{homogéneo} si los coeficientes no dependen de $\rv$.

\begin{center}
\begin{tabular}{@{}lcccc@{}}\toprule
Medio & Lineal & Dispersivo (tiempo) & Dispersión espacial & Homogéneo\\\midrule
(b) $\Pc+a\Pc^2=\epsilon_0\Ec$ & no & no & no & sí\\
(c) $a_1\ddot{\Pc}+a_2\dot{\Pc}+\Pc=\epsilon_0\chi\Ec$ & sí & sí & no & sí\\
(d) $\Pc=\epsilon_0\{a_1+a_2e^{-(x^2+y^2)}\}\Ec$ & sí & no & no & no\\\bottomrule
\end{tabular}
\end{center}
Los tres son isótropos (relación escalar, $\Pc\parallel\Ec$).

\inc{b} $\Pc+a\Pc^2=\epsilon_0\Ec$. El término $a\Pc^2$ rompe la superposición (si $\Ec_1\to\Pc_1$ y
$\Ec_2\to\Pc_2$, entonces $\Ec_1+\Ec_2$ no produce $\Pc_1+\Pc_2$ porque aparece el cruce
$2a\Pc_1\Pc_2$): es \textbf{no lineal}. No hay derivadas temporales: \textbf{no dispersivo}. No hay
derivadas espaciales: local. $a$ constante: \textbf{homogéneo}.

\emph{Desarrollo en potencias.} Proponemos $\Pc=c_1\Ec+c_2\Ec^2+c_3\Ec^3+\cdots$ y sustituimos:
\begin{equation*}
\Pc^2=c_1^2\Ec^2+2c_1c_2\Ec^3+O(\Ec^4)
\;\Rightarrow\;
c_1\Ec+(c_2+ac_1^2)\Ec^2+(c_3+2ac_1c_2)\Ec^3+O(\Ec^4)=\epsilon_0\Ec .
\end{equation*}
Igualando coeficientes de cada potencia:
\begin{align*}
\Ec^1:&\quad c_1=\epsilon_0,\\
\Ec^2:&\quad c_2=-ac_1^2=-a\epsilon_0^2,\\
\Ec^3:&\quad c_3=-2ac_1c_2=-2a\epsilon_0(-a\epsilon_0^2)=2a^2\epsilon_0^3 .
\end{align*}
(Comprobación con la solución exacta de la cuadrática:
$\Pc=\frac{-1+\sqrt{1+4a\epsilon_0\Ec}}{2a}$; con $\sqrt{1+u}=1+\frac u2-\frac{u^2}{8}+\frac{u^3}{16}-\cdots$ y
$u=4a\epsilon_0\Ec$ se obtiene $\epsilon_0\Ec-a\epsilon_0^2\Ec^2+2a^2\epsilon_0^3\Ec^3-\cdots$. Se eligió la
raíz con $+$ para que $\Pc\to0$ cuando $\Ec\to0$; la serie converge si $|4a\epsilon_0\Ec|<1$.)
Comparando con la convención de Boyd, $\Pc=\epsilon_0[\chi^{(1)}\Ec+\chi^{(2)}\Ec^2+\chi^{(3)}\Ec^3]$
(ec.~1.1.2, p.~2):
\resultado{$\chi^{(1)}=1,\qquad \chi^{(2)}=-a\epsilon_0,\qquad \chi^{(3)}=2a^2\epsilon_0^2$
\quad ($\Pc=\epsilon_0\Ec-a\epsilon_0^2\Ec^2+2a^2\epsilon_0^3\Ec^3-\cdots$). Unidades: $a$ en m$^2$/C, de
modo que $\chi^{(2)}$ queda en m/V y $\chi^{(3)}$ en m$^2$/V$^2$. Coincide con las respuestas del taller.}

\inc{c} $a_1\partial_t^2\Pc+a_2\partial_t\Pc+\Pc=\epsilon_0\chi\Ec$. Es una EDO lineal con coeficientes
constantes: \textbf{lineal}, \textbf{dispersivo} (las derivadas temporales dan memoria: $\Pc(t)$
depende de la historia de $\Ec$), local y \textbf{homogéneo}. Es exactamente el oscilador
armónico forzado amortiguado que Saleh menciona en la p.~160.

\emph{Susceptibilidad.} Con $\Ec=\mathrm{Re}\{Ee^{j\omega t}\}$, $\Pc=\mathrm{Re}\{Pe^{j\omega t}\}$, cada
$\partial_t$ se vuelve $j\omega$:
\begin{equation*}
\big(-a_1\omega^2+ja_2\omega+1\big)P=\epsilon_0\chi E
\;\Rightarrow\;
\chi(\nu)=\frac{\chi}{1-a_1\omega^2+ja_2\omega},\qquad \omega=2\pi\nu .
\end{equation*}
\emph{Comparación con Lorentz.} Dividiendo numerador y denominador de (5.5-19) (Saleh, p.~177) entre
$\nu_0^2$:
\begin{equation*}
\chi(\nu)=\chi_0\frac{\nu_0^2}{\nu_0^2-\nu^2+j\nu\Delta\nu}=\frac{\chi_0}{1-\dfrac{\nu^2}{\nu_0^2}+j\dfrac{\nu\,\Delta\nu}{\nu_0^2}} .
\end{equation*}
Identificando término a término:
$a_1(2\pi\nu)^2=\nu^2/\nu_0^2\Rightarrow\nu_0=1/(2\pi\sqrt{a_1})$; \
$a_2(2\pi\nu)=\nu\Delta\nu/\nu_0^2=\nu\Delta\nu\,4\pi^2a_1\Rightarrow\Delta\nu=a_2/(2\pi a_1)$; \ $\chi_0=\chi$.
\resultado{$\displaystyle\chi(\nu)=\frac{\chi}{1-a_1(2\pi\nu)^2+ja_2(2\pi\nu)}$;\quad
$\nu_0=\frac{1}{2\pi\sqrt{a_1}}$ (frecuencia de resonancia, $\omega_0=1/\sqrt{a_1}$),\quad
$\Delta\nu=\frac{a_2}{2\pi a_1}$ (FWHM de $\chi''$, $\sigma=a_2/a_1$),\quad $\chi_0=\chi$ (susceptibilidad de baja frecuencia).}
Físicamente: dividiendo la EDO entre $a_1$ se obtiene la forma (5.5-15),
$\ddot{\Pc}+\sigma\dot{\Pc}+\omega_0^2\Pc=\omega_0^2\epsilon_0\chi_0\Ec$, con $\omega_0^2=1/a_1$,
$\sigma=a_2/a_1$ y $\omega_0^2\chi_0=\chi/a_1$, lo que confirma $\chi_0=\chi$.

\inc{d} $\Pc=\epsilon_0\{a_1+a_2\exp[-(x^2+y^2)]\}\Ec$. Proporcional a $\Ec$: \textbf{lineal}; instantáneo:
\textbf{no dispersivo}; local; los coeficientes dependen de $x,y$: \textbf{inhomogéneo}. Es un medio
de índice gradual con $n^2(x,y)=1+a_1+a_2e^{-(x^2+y^2)}$, máximo en el eje y decreciente hacia afuera
si $a_2>0$: el perfil de una fibra o lente GRIN (Saleh, Fig.~5.2-3, p.~159). (Implícitamente $x,y$ están
medidos en alguna unidad de longitud, porque el argumento de la exponencial debe ser adimensional.)

\idea{Mire la relación constitutiva: potencias de $\Pc$ o $\Ec$ $\Rightarrow$ no lineal; derivadas en $t$
$\Rightarrow$ dispersivo; derivadas en $\rv$ $\Rightarrow$ dispersión espacial; coeficientes con $\rv$
$\Rightarrow$ inhomogéneo; tensor $\Rightarrow$ anisótropo. Al invertir $\Pc(\Ec)$ en serie, el coeficiente de
$\Ec^3$ recibe contribución del cruce $2c_1c_2$: el error típico es olvidar ese factor 2.}

%=====================================================================
\sol{3}{Onda estacionaria viajera transversal}{Saleh §5.3, ecs.~5.3-10 y 5.3-13 a 5.3-16 (p.~163); §5.4A (pp.~164--166); §8.1, ecs.~8.1-2 a 8.1-6 (pp.~292--294).}\label{sol:3}

Campo: $\E(\rv)=E_0\cos(\beta x)\,e^{-jk_zz}\,\yh$ en el vacío, $k_0=2\pi/\lambda_0=\omega/c_0$.

\inc{a} \textbf{Relación de dispersión.} Cada componente debe cumplir Helmholtz,
$\nabla^2U+k_0^2U=0$ (Saleh, ec.~5.3-16, p.~163). Con $U=E_0\cos(\beta x)e^{-jk_zz}$:
$\partial_x^2U=-\beta^2U$, $\partial_y^2U=0$, $\partial_z^2U=-k_z^2U$, de modo que
$(-\beta^2-k_z^2+k_0^2)U=0$. Además $\nabla\cdot\E=\partial_yE_y=0$, como exige (5.3-14).
\resultado{$\beta^2+k_z^2=k_0^2=(2\pi/\lambda_0)^2$. Para que la onda viaje en $z$ ($k_z$ real) hace falta
$\beta\le k_0$; si $\beta>k_0$, $k_z$ es imaginario y la onda es evanescente en $z$.}

\inc{b} \textbf{Campo magnético.} De $\nabla\times\E=-j\omega\mu_0\Hv$ (5.3-13), $\Hv=-\nabla\times\E/(j\omega\mu_0)$.
Con $\E=E_y(x,z)\yh$, las componentes del rotacional son
\begin{align*}
(\nabla\times\E)_x&=\partial_yE_z-\partial_zE_y=-\partial_zE_y=jk_zE_0\cos(\beta x)e^{-jk_zz},\\
(\nabla\times\E)_y&=\partial_zE_x-\partial_xE_z=0,\\
(\nabla\times\E)_z&=\partial_xE_y-\partial_yE_x=-\beta E_0\sin(\beta x)e^{-jk_zz}.
\end{align*}
Dividiendo entre $-j\omega\mu_0$ (y usando $1/j=-j$):
\resultado{$\displaystyle \Hv(\rv)=\frac{E_0}{\omega\mu_0}\Big[-k_z\cos(\beta x)\,\xh-j\beta\sin(\beta x)\,\zh\Big]e^{-jk_zz}$.
$H_x$ está en fase con $E_y$; $H_z$ está desfasada $90^\circ$ y desplazada un cuarto de periodo en $x$.}
Comprobación: si $\beta=0$, $\Hv=-(k_0/\omega\mu_0)E_0e^{-jk_0z}\xh=-(E_0/\eta_0)e^{-jk_0z}\xh$, la onda
plana TEM con $\E\times\Hv\parallel+\zh$ (Saleh, ec.~5.4-5, p.~165). Hay $H_z\neq0$: la onda no es TEM
respecto a $z$, es TE (el campo eléctrico sí es transversal a $z$).

\inc{c} \textbf{Vector de Poynting complejo} $\Sv=\frac12\E\times\Hv^*$ (Saleh, ec.~5.3-10, p.~163). Con
$\yh\times\xh=-\zh$ y $\yh\times\zh=\xh$:
\begin{equation*}
\Sv=\tfrac12E_y\big(H_z^*\,\xh-H_x^*\,\zh\big).
\end{equation*}
Los factores $e^{-jk_zz}$ y $e^{+jk_zz}$ se cancelan. Con
$H_x^*=-\frac{k_zE_0^*}{\omega\mu_0}\cos\beta x$ y $H_z^*=+j\frac{\beta E_0^*}{\omega\mu_0}\sin\beta x$:
\begin{align*}
S_z&=-\tfrac12E_0\cos(\beta x)\Big(-\frac{k_zE_0^*}{\omega\mu_0}\cos\beta x\Big)=\frac{k_z|E_0|^2}{2\omega\mu_0}\cos^2(\beta x),\\
S_x&=\tfrac12E_0\cos(\beta x)\Big(j\frac{\beta E_0^*}{\omega\mu_0}\sin\beta x\Big)=j\,\frac{\beta|E_0|^2}{4\omega\mu_0}\sin(2\beta x).
\end{align*}
\resultado{$\mathrm{Re}\{S_z\}=\dfrac{k_z}{k_0}\dfrac{|E_0|^2}{2\eta_0}\cos^2(\beta x)\ge0$: la potencia media fluye en $+z$,
concentrada en los vientres de $\cos\beta x$. \quad $\mathrm{Re}\{S_x\}=0$: no hay flujo neto en $x$.
\quad $\mathrm{Im}\{S_x\}=\dfrac{\beta}{2k_0}\dfrac{|E_0|^2}{2\eta_0}\sin(2\beta x)$: flujo \emph{reactivo}.}
(Se usó $\omega\mu_0=k_0\eta_0$.) La parte imaginaria del vector de Poynting complejo describe energía
que va y viene a $2\omega$ entre las regiones de campo eléctrico y magnético sin transporte neto
(el término en $e^{\pm j2\omega t}$ de la ec.~5.3-8, p.~162): en $x$ la onda es estacionaria. La
Fig.~\ref{fig:t1p03} lo muestra para $\theta=30^\circ$ ($\beta=k_0/2$).

\inc{d} \textbf{Descomposición en ondas planas.} Con $\cos\beta x=\frac12(e^{j\beta x}+e^{-j\beta x})$:
\begin{equation*}
\E=\frac{E_0}{2}\Big[e^{-j(-\beta x+k_zz)}+e^{-j(\beta x+k_zz)}\Big]\yh
=\frac{E_0}{2}\Big[e^{-j\kv_+\cdot\rv}+e^{-j\kv_-\cdot\rv}\Big]\yh,\qquad \kv_\pm=(\mp\beta,0,k_z).
\end{equation*}
Ambos tienen $|\kv_\pm|^2=\beta^2+k_z^2=k_0^2$ (por (a)) y $\E\parallel\yh\perp\kv_\pm$: cada una es una
onda plana TEM (Saleh, §5.4A, p.~165), con su $\Hv_\pm=\kv_\pm\times\E_\pm/\omega\mu_0$ (5.4-4).
\resultado{Dos ondas planas TEM de amplitud $E_0/2$ en el plano $x$--$z$, a ángulos $\pm\theta$ del eje $z$ con
$\sin\theta=\beta/k_0$, $\cos\theta=k_z/k_0$. Cada una lleva $|E_0/2|^2/2\eta_0$ en la dirección $\kv_\pm$; las
componentes $x$ de sus flujos se cancelan y las $z$ se suman, y su interferencia da la franja
$\cos^2\beta x$.}
Esto ilustra por qué $\mathrm{Re}\{S_z\}$ trae el factor $k_z/k_0=\cos\theta$: cada onda plana se
propaga oblicuamente.

\emph{Relación con la guía de espejos planos (Saleh §8.1).} Si se ponen dos espejos perfectos en
$x=\pm d/2$, la condición de frontera $E_y=0$ (tangencial) exige $\cos(\beta d/2)=0$, es decir
$\beta=m\pi/d$ con $m$ impar. Es exactamente el modo TE de orden $m$ de la guía:
$k_{xm}=m\pi/d$ (8.1-3), $\beta_m^2=k^2-m^2\pi^2/d^2$ (8.1-4, p.~293) y perfil
$\cos(m\pi x/d)$ (8.1-6, p.~294). Con $\sin(\beta x)$ en lugar de $\cos$ se obtienen los modos de $m$ par.
Así el modo guiado es la superposición de dos ondas planas que rebotan a $\pm\theta_m$, con
$\sin\theta_m=m\lambda/2d$ (8.1-2, p.~292). (Saleh usa $y$ como coordenada transversal y $E_x$; aquí los
papeles de $x$ e $y$ están intercambiados.)

\codigo{t1_p03_onda_estacionaria.py}
Salida del programa ($\lambda_0=1$ µm, $\theta=30^\circ$): $\beta=3.1416\times10^6$ m$^{-1}$,
$k_z=5.4414\times10^6$ m$^{-1}$, $\max|\mathrm{Re}\,S_x|=6\times10^{-20}$ W/m$^2$ (cero numérico),
$\max\mathrm{Re}\,S_z/(|E_0|^2/2\eta_0)=0.866=\cos30^\circ$ y $\max\mathrm{Im}\,S_x/(|E_0|^2/2\eta_0)=0.25=\beta/2k_0$.
\begin{figure}[H]\centering
\includegraphics[width=0.95\linewidth]{t1_p03_onda_estacionaria.png}
\caption{Problema 3 ($\lambda_0=1$ µm, $\theta=30^\circ$). Izquierda: instantánea del campo real
$\mathcal E_y(x,z,0)/E_0$ y flechas de $\mathrm{Re}\{\mathbf S\}$, que solo apuntan en $+z$ y son máximas en los
vientres de $\cos\beta x$. Derecha: $\mathrm{Re}\{S_z\}$ y $\mathrm{Im}\{S_x\}$ normalizados a $|E_0|^2/2\eta_0$.}
\label{fig:t1p03}
\end{figure}

\idea{La potencia media es $\mathrm{Re}\{\Sv\}$; la parte imaginaria es energía reactiva (onda estacionaria,
campo evanescente). Un campo con $\cos(\beta x)$ es siempre dos ondas planas a $\pm\theta$: así se
entienden los modos de guía (y la luz en un anillo o en una guía de SiN). Error típico: olvidar el
conjugado en $\Hv^*$ o el signo de $\yh\times\xh=-\zh$.}

%=====================================================================
\sol{4}{Campos de un láser enfocado frente al campo atómico}{Saleh ecs.~3.1-16 (p.~79), 3.1-11 (p.~77), 5.4-6 a 5.4-8 (p.~166); Boyd §1.1, ecs.~1.1.2 a 1.1.6 (pp.~2--4), ec.~2.2.16 (p.~77), §2.10 (pp.~116--122).}\label{sol:4}

Datos: $\lambda_0=1.064$ µm, $\mathcal E_p=10$ mJ, $\tau_p=10$ ns (potencia constante), $n=1.50$,
$W_0=20$ µm. Constantes: $c_0=2.998\times10^8$ m/s, $\epsilon_0=8.854\times10^{-12}$ F/m,
$\eta_0=376.7$ $\Omega$.

\inc{a} \textbf{Potencia pico.} Con potencia constante durante el pulso:
$P=\mathcal E_p/\tau_p=\SI{10e-3}{J}/\SI{10e-9}{s}=\SI{1.00e6}{W}=1$ MW.

\textbf{Intensidad pico.} Para un haz gaussiano (Saleh, ec.~3.1-16, p.~79)
$I(\rho,z)=\frac{2P}{\pi W^2(z)}e^{-2\rho^2/W^2(z)}$; en el centro del foco ($\rho=0$, $z=0$, $W=W_0$):
\begin{equation*}
I_0=\frac{2P}{\pi W_0^2}=\frac{2\times10^6\ \mathrm{W}}{\pi(20\times10^{-6}\ \mathrm m)^2}
=\frac{2\times10^6}{1.2566\times10^{-9}}\ \mathrm{W/m^2}=1.59\times10^{15}\ \mathrm{W/m^2}=1.59\times10^{11}\ \mathrm{W/cm^2}.
\end{equation*}
(El factor 2 viene de que la potencia de una gaussiana es $P=\frac12I_0\pi W_0^2$, ec.~3.1-15.)

\textbf{Fluencia pico} (energía por unidad de área en el centro): $F_0=I_0\tau_p=2\mathcal E_p/\pi W_0^2
=1.59\times10^{7}$ J/m$^2=1.59\times10^3$ J/cm$^2$.
\resultado{$P=1$ MW; \ $I_0=1.59\times10^{15}$ W/m$^2=1.59\times10^{11}$ W/cm$^2$; \ $F_0=1.59\times10^3$ J/cm$^2$.}
La fluencia está muy por encima del umbral de daño típico de un vidrio para pulsos de nanosegundos
(del orden de decenas de J/cm$^2$; valor orientativo, no del libro): en la práctica el vidrio se
dañaría. Por eso los experimentos no lineales usan pulsos cortos (ps, fs), que alcanzan
intensidades altas con poca fluencia.

\inc{b} \textbf{Campo en las dos convenciones.}
\begin{itemize}
\item \emph{Saleh} ($\Ec=\mathrm{Re}\{Ee^{j\omega t}\}$, $|E|$ = amplitud de pico del campo real):
$I=|E|^2/2\eta$, con $\eta=\eta_0/n=376.7/1.50=251.2$ $\Omega$ (ecs.~5.4-6 y 5.4-8, p.~166). Entonces
\begin{equation*}
|E_S|=\sqrt{2\eta I_0}=\sqrt{2(251.2)(1.59\times10^{15})}\ \mathrm{V/m}=\sqrt{7.99\times10^{17}}\ \mathrm{V/m}=8.94\times10^8\ \mathrm{V/m}.
\end{equation*}
\item \emph{Boyd} ($\tilde E=Ee^{-i\omega t}+\cc$, así que el campo real tiene amplitud $2|E|$):
$I=2n\epsilon_0c|E|^2$ (ec.~2.2.16, p.~77). Entonces
\begin{equation*}
|E_B|=\sqrt{\frac{I_0}{2n\epsilon_0c}}=\sqrt{\frac{1.59\times10^{15}}{2(1.50)(8.854\times10^{-12})(2.998\times10^8)}}\ \mathrm{V/m}
=\sqrt{2.00\times10^{17}}\ \mathrm{V/m}=4.47\times10^8\ \mathrm{V/m}.
\end{equation*}
\end{itemize}
\emph{Verificación de que dan la misma intensidad.} El campo real es el mismo:
$\mathrm{Re}\{E_Se^{j\omega t}\}=E_Be^{-i\omega t}+\cc=2\,\mathrm{Re}\{E_Be^{-i\omega t}\}$, luego $|E_S|=2|E_B|$.
Además $1/(2\eta)=n/(2\eta_0)=n\epsilon_0c/2$ porque $1/\eta_0=\epsilon_0c$. Así
$\frac{|E_S|^2}{2\eta}=\frac{n\epsilon_0c}{2}(2|E_B|)^2=2n\epsilon_0c|E_B|^2$: idénticas.
Numéricamente el programa da $1.5915\times10^{15}$ W/m$^2$ en ambos casos.
\resultado{$|E_S|=8.94\times10^8$ V/m (amplitud de pico); \ $|E_B|=4.47\times10^8$ V/m $=|E_S|/2$.}

\inc{c} \textbf{Campo atómico} (Boyd, p.~3), con $e=1.602\times10^{-19}$ C y $a_0=5.292\times10^{-11}$ m:
\begin{equation*}
E_{\rm at}=\frac{e}{4\pi\epsilon_0a_0^2}=\frac{(8.988\times10^9)(1.602\times10^{-19})}{(5.292\times10^{-11})^2}\ \mathrm{V/m}
=\frac{1.440\times10^{-9}}{2.800\times10^{-21}}\ \mathrm{V/m}=5.14\times10^{11}\ \mathrm{V/m}.
\end{equation*}
\textbf{Razones de polarización.} De (1.1.2) (Boyd, p.~2), $|P^{(2)}|/|P^{(1)}|=\chi^{(2)}\tilde E^2/\chi^{(1)}\tilde E
=(\chi^{(2)}/\chi^{(1)})\tilde E\sim\tilde E/E_{\rm at}$ y $|P^{(3)}|/|P^{(1)}|\sim(\tilde E/E_{\rm at})^2$ (ecs.~1.1.3--1.1.4,
p.~3). Aquí $\tilde E$ es el \emph{campo real} que entra en la serie de potencias, cuya amplitud es
$|E_S|=2|E_B|=8.94\times10^8$ V/m (usar $|E_B|$ daría razones dos veces menores: error de convención típico):
\begin{equation*}
\frac{|P^{(2)}|}{|P^{(1)}|}\sim\frac{8.94\times10^8}{5.14\times10^{11}}=1.7\times10^{-3},\qquad
\frac{|P^{(3)}|}{|P^{(1)}|}\sim(1.74\times10^{-3})^2=3.0\times10^{-6}.
\end{equation*}
\resultado{$E_{\rm at}=5.14\times10^{11}$ V/m; \ $|P^{(2)}|/|P^{(1)}|\sim1.7\times10^{-3}$; \ $|P^{(3)}|/|P^{(1)}|\sim3\times10^{-6}$.}
Aunque son razones pequeñas, la polarización no lineal radia coherentemente: si hay ajuste de fases,
los campos de todos los dipolos se suman en amplitud a lo largo de muchos miles de longitudes de onda
y la señal generada crece como $L^2$ (Boyd, §2.1--2.2). Por eso un $10^{-3}$ en la fuente basta para
convertir una fracción apreciable de la potencia. Para que $P^{(2)}\sim P^{(1)}$ se necesitaría
$I_{\rm at}\approx3.5\times10^{16}$ W/cm$^2$ (Boyd, ec.~1.1.6, p.~4), cinco órdenes por encima.

\inc{d} \textbf{Rango de Rayleigh en el vidrio.} De (3.1-11), $z_0=\pi W_0^2/\lambda$ con $\lambda=\lambda_0/n$ la
longitud de onda \emph{en el medio} (Saleh, p.~77):
\begin{equation*}
z_0=\frac{\pi W_0^2n}{\lambda_0}=\frac{\pi(20\times10^{-6})^2(1.50)}{1.064\times10^{-6}}\ \mathrm m=1.77\times10^{-3}\ \mathrm m=1.77\ \mathrm{mm},
\qquad 2z_0=3.54\ \mathrm{mm}.
\end{equation*}
\resultado{$z_0=1.77$ mm (profundidad de foco $2z_0=3.54$ mm; divergencia en el vidrio
$\theta_0=\lambda_0/(n\pi W_0)=11.3$ mrad).}
\textbf{Por qué importa.} La intensidad solo es alta en una longitud $\approx2z_0$ (el parámetro
confocal $b=2z_0$ de Boyd, §2.10.2, p.~117); fuera de ella cae como $1/[1+(z/z_0)^2]$ (Fig.~\ref{fig:t1p04}).
En generación de armónicos hay un compromiso: enfocar más fuerte (menor $W_0$) aumenta $I_0\propto1/W_0^2$,
pero acorta la región útil ($z_0\propto W_0^2$). Además, al cruzar el foco el haz adquiere la fase de
Gouy ($\pi$ en total), que se suma al desajuste de fases $\Delta k$: con enfoque fuerte la THG en un
medio con $\Delta k\le0$ se cancela, y el óptimo para SHG es la cintura en el centro del cristal con
$L/b=2.84$ y $\Delta k=3.2/L$ (Boyd--Kleinman; Boyd, §2.10.3, pp.~119--122). Con $b=3.54$ mm, el
cristal óptimo mediría unos 10 mm.

\codigo{t1_p04_foco_laser.py}
\begin{figure}[H]\centering
\includegraphics[width=0.85\linewidth]{t1_p04_foco.png}
\caption{Problema 4: intensidad normalizada $I(\rho,z)/I_0$ alrededor del foco dentro del vidrio
($W_0=20$ µm, $n=1.50$, $\lambda_0=1.064$ µm), con $\pm W(z)$ y la profundidad de foco $2z_0=3.54$ mm.}
\label{fig:t1p04}
\end{figure}

\idea{Cadena de cálculo: energía/duración $\to P$; $I_0=2P/\pi W_0^2$ (¡el 2!); $|E|$ con la convención que
toque. Saleh: $I=|E|^2/2\eta=n\epsilon_0c|E|^2/2$; Boyd: $I=2n\epsilon_0c|E|^2$; por tanto $|E_S|=2|E_B|$. Para
estimar $P^{(2)}/P^{(1)}\sim E/E_{\rm at}$ use la amplitud del campo real. Dentro de un medio, $z_0=\pi W_0^2n/\lambda_0$
(se olvida la $n$ con frecuencia).}

%=====================================================================
\sol{5}{Ondas esféricas y cilíndricas}{Saleh §2.2A--B, ecs.~2.2-7, 2.2-10 y 2.2-15 (pp.~43--46); Probs.~2.2-3 a 2.2-5 (p.~73).}\label{sol:5}

\inc{a} Para una función que solo depende de $r$, el laplaciano en coordenadas esféricas es
$\nabla^2U=\frac1{r^2}\frac{d}{dr}\Big(r^2\frac{dU}{dr}\Big)$. Con $U=(A_0/r)e^{-jkr}$ (Saleh, ec.~2.2-15, p.~46):
\begin{align*}
\frac{dU}{dr}&=A_0\Big(-\frac1{r^2}-\frac{jk}{r}\Big)e^{-jkr},\qquad
r^2\frac{dU}{dr}=-A_0(1+jkr)e^{-jkr},\\
\frac{d}{dr}\Big(r^2\frac{dU}{dr}\Big)&=-A_0\big[jk-jk(1+jkr)\big]e^{-jkr}=-A_0\big[jk-jk+k^2r\big]e^{-jkr}=-A_0k^2r\,e^{-jkr}.
\end{align*}
Dividiendo entre $r^2$: $\nabla^2U=-k^2\frac{A_0}{r}e^{-jkr}=-k^2U$, es decir:
\resultado{$\nabla^2U+k^2U=0$ para $r\neq0$ (Helmholtz, ec.~2.2-7, p.~43).}
(Atajo útil: $\nabla^2U=\frac1r\frac{d^2(rU)}{dr^2}$, y $rU=A_0e^{-jkr}$ da $-k^2rU$ de inmediato. En $r=0$ la
función es singular: $\nabla^2(1/r)=-4\pi\delta(\rv)$ representa la fuente puntual.)

\inc{b} En óptica escalar $I=|U|^2$ (Saleh, ec.~2.2-10, p.~43), así que $I(r)=|A_0|^2/r^2$. La potencia que
atraviesa una esfera de radio $r$ es $P=\oint I\,dA=I(r)\,4\pi r^2=4\pi|A_0|^2$, independiente de $r$
(conservación de la energía). Despejando:
\begin{equation*}
I(r)=\frac{P}{4\pi r^2}=\frac{100\ \mathrm W}{4\pi(1\ \mathrm m)^2}=7.96\ \mathrm{W/m^2}=7.96\times10^{-4}\ \mathrm{W/cm^2}.
\end{equation*}
\resultado{$I=P/4\pi r^2=7.96$ W/m$^2$ en $r=1$ m para $P=100$ W.}
Compárese con el Prob.~4: una potencia $10^4$ veces mayor (1 MW), concentrada en una cintura de 20 µm, da una
intensidad $2\times10^{14}$ veces mayor. La óptica no lineal necesita confinamiento, no solo potencia.

\inc{c} Frentes cilíndricos centrados en el eje $y$: $U=U(\rho)$ con $\rho=\sqrt{x^2+z^2}$, sin dependencia en $y$ ni
en el ángulo azimutal. El laplaciano en coordenadas cilíndricas (eje $y$) se reduce a
$\nabla^2U=\frac1\rho\frac{d}{d\rho}\big(\rho\frac{dU}{d\rho}\big)$, y Helmholtz queda
$U''+\frac1\rho U'+k^2U=0$, la ecuación de Bessel de orden cero. La solución exacta saliente (con $e^{j\omega t}$) es
la función de Hankel $H_0^{(2)}(k\rho)\approx\sqrt{2/\pi k\rho}\,e^{-j(k\rho-\pi/4)}$ para $k\rho\gg1$.

Sin recurrir a Bessel, probamos el ansatz $U=A\rho^{-1/2}e^{-jk\rho}$:
\begin{align*}
U'&=A\Big(-\tfrac12\rho^{-3/2}-jk\rho^{-1/2}\Big)e^{-jk\rho},\qquad
\rho U'=A\Big(-\tfrac12\rho^{-1/2}-jk\rho^{1/2}\Big)e^{-jk\rho},\\
(\rho U')'&=A\Big[\tfrac14\rho^{-3/2}-\tfrac{jk}{2}\rho^{-1/2}-jk\Big(-\tfrac12\rho^{-1/2}-jk\rho^{1/2}\Big)\Big]e^{-jk\rho}
=A\Big[\tfrac14\rho^{-3/2}-k^2\rho^{1/2}\Big]e^{-jk\rho}.
\end{align*}
Dividiendo entre $\rho$: $\nabla^2U=-k^2U+\frac{1}{4\rho^2}U$, de modo que $\nabla^2U+k^2U=\frac{U}{4\rho^2}$. El residuo
relativo es $1/(4k^2\rho^2)$, despreciable si $k\rho\gg1$ (a unas pocas longitudes de onda del eje).
\resultado{$U(\rho)\approx\dfrac{A}{\sqrt\rho}e^{-jk\rho}$ $(k\rho\gg1)$; \quad $I=|U|^2=\dfrac{|A|^2}{\rho}\propto\dfrac1\rho$.}
\textbf{Por qué $1/\rho$.} La potencia por unidad de longitud en $y$ que atraviesa un cilindro de radio $\rho$ es
$P_\ell=I\cdot2\pi\rho=2\pi|A|^2$, constante. Como el área del cilindro crece solo linealmente con $\rho$ (y no
como $r^2$ en la esfera), la intensidad cae como $1/\rho$ y la amplitud como $1/\sqrt\rho$. Es el campo lejano de
una rendija larga o de una fuente lineal.

\idea{Conservación de la energía fija el decaimiento: esfera $I\propto1/r^2$ (amplitud $1/r$), cilindro
$I\propto1/\rho$ (amplitud $1/\sqrt\rho$), plano: constante. Use $\nabla^2U=\frac1r(rU)''$ en esféricas. En esta
sección Saleh usa $I=|U|^2$ (óptica escalar), no $|E|^2/2\eta$.}

%=====================================================================
\sol{6}{Helmholtz paraxial y el haz gaussiano}{Saleh §2.2C, ecs.~2.2-20 a 2.2-23 (pp.~47--48); §3.1A, ecs.~3.1-4 a 3.1-11 (pp.~76--77); Ejercicio~3.1-2 (p.~84).}\label{sol:6}

\inc{a} Sustituimos $U=A(\rv)e^{-jkz}$ (2.2-20) en $\nabla^2U+k^2U=0$. Las derivadas transversales solo actúan
sobre $A$: $\nabla_T^2U=(\nabla_T^2A)e^{-jkz}$. En $z$, con la regla del producto:
\begin{align*}
\partial_zU&=(\partial_zA-jkA)e^{-jkz},\\
\partial_z^2U&=(\partial_z^2A-jk\partial_zA-jk\partial_zA+(-jk)^2A)e^{-jkz}=(\partial_z^2A-2jk\partial_zA-k^2A)e^{-jkz}.
\end{align*}
Sumando y añadiendo $k^2U$, el término $-k^2A$ se cancela con $+k^2A$:
\begin{equation*}
\big[\nabla_T^2A+\partial_z^2A-2jk\,\partial_zA\big]e^{-jkz}=0 .
\end{equation*}
\emph{Aproximación.} (2.2-21) dice que en una distancia $\Delta z=\lambda$ el cambio de $A$ es pequeño:
$|\partial_zA|\lambda\ll|A|$, o sea $|\partial_zA|\ll k|A|$. Aplicando lo mismo a la función $\partial_zA$ (que también
varía lentamente) se tiene $|\partial_z^2A|\ll k|\partial_zA|$, que implica (2.2-22). Por tanto
$\partial_z^2A$ es despreciable frente a $2k\partial_zA$ y queda
\resultado{$\nabla_T^2A-j2k\,\partial_zA=0$ \quad (ecuación paraxial de Helmholtz, Saleh ec.~2.2-23, p.~48).}
Ojo: \emph{no} se desprecia $\nabla_T^2A$; la variación transversal puede ser rápida (en una escala $W$), solo la
longitudinal debe ser lenta. Es la misma SVEA que usa Boyd para las ecuaciones de ondas acopladas.

\inc{b} Sea $A=\frac{A_1}{q}f$, $f=\exp\!\big(-\frac{jk\rho^2}{2q}\big)$, $q=q(z)$, $q'=dq/dz$ (Saleh, ec.~3.1-4, p.~76).

\emph{Laplaciano transversal.} $\partial_xf=-\frac{jkx}{q}f$, \
$\partial_x^2f=\big[-\frac{jk}{q}+\big(\frac{-jkx}{q}\big)^2\big]f=\big[-\frac{jk}{q}-\frac{k^2x^2}{q^2}\big]f$. Igual en $y$; sumando,
\begin{equation*}
\nabla_T^2A=\frac{A_1}{q}f\Big[-\frac{2jk}{q}-\frac{k^2\rho^2}{q^2}\Big].
\end{equation*}
\emph{Derivada en $z$.} $\partial_z(1/q)=-q'/q^2$ y $\partial_zf=f\cdot\frac{jk\rho^2}{2q^2}q'$, luego
\begin{equation*}
\partial_zA=\frac{A_1}{q}f\Big[-\frac{q'}{q}+\frac{jk\rho^2q'}{2q^2}\Big]
\;\Rightarrow\;
-2jk\,\partial_zA=\frac{A_1}{q}f\Big[\frac{2jkq'}{q}+\frac{k^2\rho^2q'}{q^2}\Big]
\end{equation*}
(se usó $(-2jk)(jk)/2=k^2$). \emph{Suma.}
\begin{equation*}
\nabla_T^2A-2jk\,\partial_zA=\frac{A_1}{q}f\,(q'-1)\Big[\frac{2jk}{q}+\frac{k^2\rho^2}{q^2}\Big].
\end{equation*}
El corchete no es idénticamente cero (depende de $\rho$), así que la ecuación se cumple para todo $\rho$ y $z$
\emph{si y solo si} $q'=1$. Integrando: $q(z)=z+q_0$ con $q_0$ constante (compleja). Si $q_0$ fuera real
tendríamos la onda paraboloidal centrada en $z=-q_0$; para que el haz tenga un perfil transversal acotado
se necesita $\mathrm{Im}\{1/q\}<0$, y se elige $q_0=jz_0$ con $z_0>0$ real:
\resultado{$q(z)=z+jz_0$ (Saleh, ec.~3.1-5, p.~76).}
\emph{Parámetros del haz.} Racionalizando,
$\dfrac1q=\dfrac{z-jz_0}{z^2+z_0^2}$. Igualando con $\dfrac1q=\dfrac1R-j\dfrac{\lambda}{\pi W^2}$ (3.1-6):
\begin{align*}
\text{Parte real:}&\quad R(z)=\frac{z^2+z_0^2}{z}=z\Big[1+\Big(\frac{z_0}{z}\Big)^2\Big],\\
\text{Parte imaginaria:}&\quad \frac{\lambda}{\pi W^2}=\frac{z_0}{z^2+z_0^2}\;\Rightarrow\;
W^2(z)=\frac{\lambda z_0}{\pi}\Big[1+\Big(\frac{z}{z_0}\Big)^2\Big],\\
&\quad\Rightarrow\;W(z)=W_0\sqrt{1+(z/z_0)^2},\ \ W_0=\sqrt{\lambda z_0/\pi}.
\end{align*}
Para la fase de Gouy escribimos el prefactor con $A_0=A_1/jz_0$ (Saleh, p.~77):
\begin{equation*}
\frac{A_1}{q}=\frac{jz_0A_0}{z+jz_0}=\frac{A_0}{1-jz/z_0}=\frac{A_0}{\sqrt{1+(z/z_0)^2}\,e^{-j\zeta}}=A_0\frac{W_0}{W(z)}e^{j\zeta(z)},
\qquad \zeta(z)=\tan^{-1}\frac{z}{z_0}.
\end{equation*}
Y el exponente: $-\frac{jk\rho^2}{2q}=-\frac{jk\rho^2}{2R}-\frac{k\lambda\rho^2}{2\pi W^2}=-\frac{jk\rho^2}{2R}-\frac{\rho^2}{W^2}$ (pues $k\lambda=2\pi$).
\resultado{$W(z)=W_0\sqrt{1+(z/z_0)^2}$, \ $R(z)=z[1+(z_0/z)^2]$, \ $\zeta(z)=\tan^{-1}(z/z_0)$, \ $W_0=\sqrt{\lambda z_0/\pi}$
(ecs.~3.1-8 a 3.1-11, p.~77); juntando todo se obtiene la amplitud (3.1-7).}

\inc{c} \textbf{Validez paraxial (Ejercicio~3.1-2).} Hay que comparar $|\partial_zA|$ con $k|A|$. De (b),
\begin{equation*}
\frac{\partial_zA}{A}=-\frac1q+\frac{jk\rho^2}{2q^2}.
\end{equation*}
El primer término vale a lo sumo $1/|q|\le1/z_0$. El segundo, evaluado en el borde del haz $\rho=W(z)$, vale
$\frac{kW^2}{2|q|^2}=\frac{k}{2}\frac{\lambda z_0}{\pi}\frac{1+(z/z_0)^2}{z_0^2[1+(z/z_0)^2]}=\frac{1}{z_0}$ (usando $|q|^2=z^2+z_0^2$ y $k\lambda=2\pi$).
Así $|\partial_zA|/|A|\sim1/z_0$, y la condición (2.2-21) es
\begin{equation*}
\frac{|\partial_zA|}{k|A|}\sim\frac{1}{kz_0}=\frac{\lambda}{2\pi z_0}=\frac{\lambda^2}{2\pi^2W_0^2}=\frac{\theta_0^2}{2}\ll1,
\qquad \theta_0=\frac{\lambda}{\pi W_0}.
\end{equation*}
Es decir, la aproximación paraxial vale si la divergencia es pequeña, $\theta_0\ll1$. Algunos valores:
\begin{center}
\begin{tabular}{@{}cccc@{}}\toprule
$W_0/\lambda$ & $\theta_0$ (rad) & $\theta_0$ (grados) & $\theta_0^2/2$\\\midrule
0.5 & 0.64 & 36 & 0.20\\
1 & 0.32 & 18 & 0.051\\
2 & 0.16 & 9.1 & 0.013\\
5 & 0.064 & 3.6 & 0.0020\\\bottomrule
\end{tabular}
\end{center}
\resultado{La aproximación es razonable para $W_0\gtrsim\lambda$ (error relativo $\lesssim5$\,\%, $\theta_0\lesssim0.3$ rad) y falla
para $W_0\lesssim\lambda/2$ ($\theta_0\gtrsim0.6$ rad, error $\sim20$\,\%).}
Para guías de SiN con modos de $\sim1$ µm a 1.55 µm se está en el límite: por eso los modos de guías de alto
contraste se calculan con solvers vectoriales completos (FDE) y no con la teoría paraxial.

\idea{La ecuación paraxial sale de despreciar $\partial_z^2A$ frente a $2k\partial_zA$; nada más. El haz gaussiano es
la onda paraboloidal con $z\to z+jz_0$: toda la física está en $1/q=1/R-j\lambda/\pi W^2$. La validez paraxial se
resume en $1/kz_0=\theta_0^2/2\ll1$. Errores típicos: el signo de $\zeta$ (con $e^{j\omega t}$ aparece como $+j\zeta$ en la
amplitud) y olvidar $\lambda=\lambda_0/n$ en un medio.}

%=====================================================================
\sol{7}{Pulso con chirp en el vacío}{Saleh §5.1 (ecs.~5.1-2 y 5.1-13, pp.~152--155), §5.4A (ecs.~5.4-5 a 5.4-8, pp.~165--166); frecuencia instantánea y chirp: §22.1A, ecs.~22.1-4 y 22.1-5 (pp.~939--940).}\label{sol:7}

Campo: $\Ec(z,t)=\mathrm{Re}\{f(t-z/c_0)\}\,\xh$, con $f(t)=e^{-t^2/\tau^2}e^{j2\pi\nu_0t+jbt^2}$. El enunciado no da
amplitud; escribimos $E_0f$ con $E_0$ real en V/m (en el enunciado, $E_0=1$ V/m).

\inc{a} \textbf{Descripción.}
\begin{itemize}
\item \emph{Polarización}: lineal en $\xh$ (el vector campo siempre apunta en $\pm\xh$).
\item \emph{Dirección}: depende de $z$ y $t$ solo a través de $t-z/c_0$; es una onda que viaja en $+z$ a $c_0$ sin
deformarse (el vacío no es dispersivo). Es una onda plana pulsada (sin dependencia en $x,y$).
\item \emph{Envolvente}: gaussiana $e^{-(t-z/c_0)^2/\tau^2}$; el campo cae a $1/e$ en $|t-z/c_0|=\tau$ y la intensidad
($\propto$ envolvente$^2$) a $1/e^2$; su FWHM en intensidad es $\tau\sqrt{2\ln2}=1.18\tau$.
\item \emph{Portadora}: $\cos[2\pi\nu_0(t-z/c_0)+b(t-z/c_0)^2]$: frecuencia central $\nu_0$ con un \emph{chirp} lineal
(la fase cuadrática $bt^2$ hace que la frecuencia cambie a lo largo del pulso).
\end{itemize}

\inc{b} \textbf{Campo magnético.} Como el campo no es monocromático, usamos Maxwell en el dominio del tiempo,
$\nabla\times\Ec=-\mu_0\partial_t\Hc$ (5.1-2). Sea $u=t-z/c_0$ y $\mathcal E_x(z,t)=E_0\,\mathrm{Re}\{f(u)\}$. Solo hay
$\mathcal E_x(z,t)$, así que $\nabla\times\Ec=\partial_z\mathcal E_x\,\yh$, y por la regla de la cadena
$\partial_z=-\frac1{c_0}\frac{d}{du}$:
\begin{equation*}
\nabla\times\Ec=-\frac{E_0}{c_0}\mathrm{Re}\{f'(u)\}\,\yh=-\mu_0\frac{\partial\Hc}{\partial t}.
\end{equation*}
Como $\partial_t=\frac{d}{du}$, se integra en $u$: $\mathcal H_y=\frac{E_0}{\mu_0c_0}\mathrm{Re}\{f(u)\}$ (la constante de
integración es cero porque el campo se anula lejos del pulso). Con $\mu_0c_0=\eta_0$:
\resultado{$\Hc(z,t)=\dfrac{E_0}{\eta_0}\mathrm{Re}\{f(t-z/c_0)\}\,\yh$, \ $\eta_0=\sqrt{\mu_0/\epsilon_0}\approx377$ $\Omega$.}
Las demás ecuaciones se cumplen: $\nabla\cdot\Ec=\partial_x\mathcal E_x=0$, $\nabla\cdot\Hc=0$, y
$\nabla\times\Hc=-\partial_z\mathcal H_y\,\xh=\frac{E_0}{\eta_0c_0}\mathrm{Re}\{f'\}\xh=\epsilon_0\partial_t\Ec$ porque
$1/(\eta_0c_0)=\epsilon_0$. Igual que en la onda plana monocromática (5.4-5): $\Ec\perp\Hc\perp\zh$,
$|\Ec|/|\Hc|=\eta_0$ instante a instante, y $\Ec\times\Hc\parallel+\zh$.

\inc{c} \textbf{Frecuencia instantánea.} La fase total de $f$ es $\phi(t)=2\pi\nu_0t+bt^2$, así que
\begin{equation*}
\nu_i(t)=\frac{1}{2\pi}\frac{d\phi}{dt}=\frac{1}{2\pi}(2\pi\nu_0+2bt)=\nu_0+\frac{b\,t}{\pi}.
\end{equation*}
Equivale a (22.1-4) de Saleh (p.~939) con fase de envolvente $\varphi(t)=bt^2$.
\resultado{$\nu_i(t)=\nu_0+bt/\pi$: con $b>0$ la frecuencia crece con el tiempo, \emph{up-chirp} (el frente del pulso es
rojo y la cola azul).}
En la notación de Saleh, $\varphi=at^2/\tau^2$, el parámetro de chirp es $a=\tfrac12\varphi''\tau^2=b\tau^2$ (22.1-5, p.~940);
$a>0$ es up-chirp, consistente. \emph{Advertencia de convención:} con $e^{-i\omega t}$ (Boyd, Agrawal) el mismo campo real
se escribe con fase $-bt^2$ en la envolvente, y el signo que indica up-chirp se invierte. Lo que no cambia es la física:
$\nu_i$ crece en el tiempo.

\inc{d} \textbf{Números.} $\tau=20$ fs, $\nu_0=375$ THz ($\lambda_0=799$ nm, periodo óptico 2.67 fs), $b=5\times10^{27}$ s$^{-2}$:
\begin{equation*}
\Delta\nu_i=\nu_i(\tau)-\nu_i(-\tau)=\frac{2b\tau}{\pi}=\frac{2(5\times10^{27}\ \mathrm{s^{-2}})(20\times10^{-15}\ \mathrm s)}{\pi}
=\frac{2\times10^{14}}{\pi}\ \mathrm{Hz}=63.7\ \mathrm{THz}.
\end{equation*}
Es decir, $\nu_i$ va de 343.2 THz (874 nm) en $t=-\tau$ a 406.8 THz (737 nm) en $t=+\tau$: un 17\,\% de $\nu_0$.
\resultado{$\Delta\nu_i=2b\tau/\pi=63.7$ THz entre $t=-\tau$ y $t=+\tau$.}
Orden de magnitud: el parámetro de chirp es $a=b\tau^2=2$, así que el espectro es $\sqrt{1+a^2}=2.24$ veces más ancho que
el de un pulso sin chirp de la misma duración (FWHM espectral de 18.7 THz sin chirp contra 41.9 THz con chirp). Ese
exceso de ancho de banda es el que permite comprimir el pulso en un medio con la GVD de signo adecuado (Problema~17).

\inc{e} \textbf{Intensidad.} Para una onda plana en el vacío, $I=|E|^2/2\eta_0$ (5.4-8) con $E$ la envolvente compleja,
promediando en un ciclo óptico (válido porque $\tau=20$ fs $\gg$ 2.67 fs). Aquí $|E_0f|^2=E_0^2e^{-2(t-z/c_0)^2/\tau^2}$; la
fase (incluido el chirp) desaparece al tomar el módulo:
\resultado{$I(z,t)=\dfrac{E_0^2}{2\eta_0}\exp\!\Big[-\dfrac{2(t-z/c_0)^2}{\tau^2}\Big]$ \ (con $E_0=1$ V/m, $I_{\rm max}=1.33\times10^{-3}$ W/m$^2$).}
Equivale a $\langle\Ec\times\Hc\rangle=\frac{E_0^2}{\eta_0}\langle\mathrm{Re}^2\{f\}\rangle\zh$, con
$\langle\cos^2\rangle=\frac12$. El chirp no cambia $I(t)$, solo el espectro: un medidor de intensidad lento no lo distingue
de un pulso sin chirp.

\codigo{t1_p07_pulso_chirp.py}
\begin{figure}[H]\centering
\includegraphics[width=0.8\linewidth]{t1_p07_pulso_chirp.png}
\caption{Problema 7. Arriba: campo real $\mathcal E_x(0,t)/E_0$, envolvente $\pm e^{-t^2/\tau^2}$ e intensidad
$e^{-2t^2/\tau^2}$; los ciclos se aprietan hacia la derecha (up-chirp). Abajo: frecuencia instantánea
$\nu_i=\nu_0+bt/\pi$ (THz) frente a $t-z/c_0$ (fs), con $\Delta\nu_i=63.7$ THz entre $\pm\tau$.}
\label{fig:t1p07}
\end{figure}

\idea{$\nu_i=\frac1{2\pi}d\phi/dt$ con la fase \emph{total}; una fase cuadrática $bt^2$ da $\nu_i$ lineal en $t$ con pendiente
$b/\pi$ (no $b/2\pi$: el 2 de derivar $t^2$ se cancela con el de $2\pi$). Con $e^{j\omega t}$, $b>0$ es up-chirp; con
$e^{-i\omega t}$ el signo se invierte. En el vacío $\Hc=\zh\times\Ec/\eta_0$ también para pulsos, y la intensidad solo ve la
envolvente.}
